\subsection[Merge Sort]{Merge Sort \cite{lobo2020performance}}
    \subsubsection{Overview}
        \begin{itemize}
            \item \textbf{History and Origin:} Merge Sort is a divide-and-conquer algorithm that was invented by John von Neumann in 1945.
            \item \textbf{Efficiency:} Consistently performs well on large datasets with a time complexity of O(n log n).
            \item \textbf{Application:} Commonly used in external sorting and when stable sorting is required.
        \end{itemize}
    \subsubsection{Idea}
        \begin{itemize}
            \item The array is recursively divided into two halves until each sub-array contains a single element.
            \item The sub-arrays are then merged back together in a sorted manner.
        \end{itemize}
    \subsubsection{Step-by-Step Algorithm}
        \begin{itemize}
            \item Consider an array of n elements, a[0\ldots (n - 1)].
            \item \textbf{Step 1:} Start an increase variable \textit{i} that is from 1 to n - 1.
            \item \textbf{Step 2:} Find the correct position \textit{pos} in a[0\ldots (i - 1)] to insert a[i], where a[pos - 1] < a[i] < a[pos].
            \begin{itemize}
                \item Set val = a[i], then move forward a[pos\ldots (i - 1)] to one element. After that, set a[pos] = val.
                \item Increase \textit{i} by 1.
            \end{itemize}
            \item \textbf{Step 3:} If i < n return to \textbf{Step 2}. Otherwise, stop the algorithm. 
        \end{itemize}

    \subsubsection{Pseudo Code}
        \begin{itemize}
            \item \textbf{Algorithm:} Insertion Sort
            \item \textbf{Input:} An array A of n elements
        \end{itemize}
        \begin{verbatim}
        1. For i = 1 to n - 1:
        2.    val = A[i]
        3.    j = i - 1
        4.    While j >= 0 and A[j] > key:
        5.       A[j + 1] = A[j]
        6.       j = j - 1
        7.    A[j + 1] = key
        \end{verbatim}
    \subsubsection{Complexity}
        \begin{itemize}
            \item \textbf{Time Complexity:}
            \begin{itemize}
                \item Best case: $O(n)$
                    \subitem When the input is already sorted.
                \item Average case: $O(n^2)$
                    \subitem When the input is random.
                \item Worst case: $O(n^2)$ 
                    \subitem When the input is reverse-sorted.
            \end{itemize}
            \item \textbf{Space Complexity:} O(1) 
                \subitem In-place algorithm. Requires only a constant amount of extra memory.
        \end{itemize}
        
    \subsubsection{Variants and Optimizations}
        \begin{itemize}
            \item \textbf{The "Hit" Method (Enhanced Insertion Sort)\cite{sodhi2013enhanced}:}
            \begin{itemize}
                \item \textbf{Problem:} Conventional Insertion Sort is inefficient for reverse-sorted lists due to maximum comparisons.
                \item \textbf{Solution:} Immediately compare $A[i]$ with $A[0]$ (the current minimum).
                \item \textbf{Benefit:} If $A[i] < A[0]$, it is moved to the front instantly, bypassing all intermediate comparisons and significantly reducing execution time.
            \end{itemize}
            \item \textbf{Binary Insertion Sort:}
            \begin{itemize}
                \item Use Binary Search to find the position of an element within the sorted sub-array which is greater than or equal to A[i].
                \item While Binary Search reduces the number of comparisons to $O(n \log n)$, the complexity of data movement remains $O(n^2)$ because elements must still be shifted to make space for the new entry.
            \end{itemize}
        \end{itemize}